% This sets the page numbers to be printed in lowercase roman numerals, starts
% the page numbering from i, and prohibits any numbering of sectional divisions.
\frontmatter

% Generate the title page.
\titlepage[crests/officialmono]

\tableofcontents
\clearforchapter\listoffigures
\clearforchapter\listoftables

\frontchapter{Acknowledgements}

I would like to thank my dissertation supervisor, Sean, for his guidance, feedback and support during this project. His suggestions helped me clarify the direction of the work and improve both the project implementation and the dissertation writing. I am also grateful to Thomas, the programme leader, for his support during my MSc study, and to Professor Kurt for his advice in the dissertation module and related guidance. I would also like to thank my family and friends for their encouragement and support during the development of the project and the writing of this dissertation.

\frontchapter{Declaration}

This dissertation is submitted to the University of Warwick as material for my MSc Games Engineering degree. The dissertation was completed by myself and has not been submitted for any other degree. Except for the materials that are clearly cited, the software, test data and analysis described in this dissertation all come from this project.

\frontchapter{Abstract}

This project developed a visual behaviour tree plugin for Godot 4.6 and placed the main research focus on the display problems of large behaviour trees. When a traditional behaviour tree editor grows larger, it occupies a large canvas, and dragging, connections and unrelated branches can also interfere with observation at the same time. To deal with these problems, the current project implemented five display optimisation functions: Smart Drag Reflow, Adaptive Zoom Detail, Readable Edge Overlay, Related Node Focus and Fisheye Focus.

To verify the influence of these functions on developers in actual development, the project used five real behaviour trees with different node counts to build a game with a certain level of complexity, and tested it under three common screen sizes. By testing behaviour trees with different node counts and screens with different sizes, the project compared the actual optimisation effect of each improved function. It also introduced the project developer's structured usage evaluation of each function.

The results show that Adaptive Zoom Detail has a relatively stable effect across different tree sizes and screen sizes, so it is the most suitable general display optimisation function. Smart Drag Reflow can handle overlap caused by node movement, and Related Node Focus can highlight the structural range of selected nodes. Fisheye Focus is more obvious on small screens and large behaviour trees, but stronger local magnification also increases some node overlap. Readable Edge Overlay can improve line segments blocked by node cards, but it mainly applies to local connection problems.

All function tests met the preset success conditions, and did not change the behaviour tree structure, saved coordinates or sibling execution order. The experiment shows that these functions can improve some measurable display and editing states of large behaviour trees, and also shows that different functions are suitable for different development environments.

Keywords: behaviour tree; Godot; node-link editor; semantic zoom; focus and context; graph layout; game artificial intelligence

\mainmatter

